% Canonical standalone Question 2 section.  The surrounding document_en.tex is
% an integration draft; this file is the Q2-only deliverable.
\section{Question 2: Capacity-Valid Scheduling and Method Selection}

\subsection{Model and policy definitions}
Each GPU model is evaluated as its own pool. Job $j$ has release time $r_j$,
duration $p_j$, $w_j$ workers, per-worker GPU request $g_j$, and CPU request
$c_j$. It starts at $s_j\ge r_j$, finishes at $s_j+p_j$, and all workers
start together. A placement is feasible only when, on every selected node $i$,
\[
g_jx_{j,i}\le G_i^{\rm free}(t),\qquad
c_jx_{j,i}\le C_i^{\rm free}(t),\qquad \sum_i x_{j,i}=w_j .
\]
The simulator is non-preemptive, uses exact integer-second durations, and
commits a gang placement atomically. All comparisons use the same numeric
First-Fit-by-ID (FF-id), raw CSV, loader ordering, and explicit
\texttt{shared-timer} event mode.

Ordinary \texttt{SPT-HOL} orders arrived jobs by
$(p_j,r_j,\mathrm{job\_id})$. It has no aging threshold, urgent promotion, or
reservation; if its head is infeasible, that dispatch epoch stops. Ordinary
\texttt{SPT-FILL} uses the same order but may inspect later feasible jobs. Its
\(K=200\) is a per-pool, per-timestamp budget only after a blocking prefix has
been encountered: feasible prefix checks do not consume it, successful and
failed bypass attempts both consume it, retries may consume it again, and a
successful start does not reset the timestamp budget. Shared timer ticks give
additional dispatch opportunities and can therefore change the observed
budgeted fill behavior. These statements describe only the tested
\(K=200\), FF-id, shared-timer configuration; they do not establish behavior
for an unbounded FILL policy. Ordinary \texttt{SPT-HOL} has no aging protection.

The formal A10 runs use $\theta=604800$ seconds (168 hours), $K=200$,
\texttt{shared-timer}, $H=7,776,000$ seconds, and complete drain. The threshold
is exploratory, not an optimum, and is an experiment configuration separate
from ordinary \texttt{SPT-HOL}. Protected urgent policies are reported only
for the verified 65,536-job A10 prefix.

\subsection{Job-count lower bound and exact waiting accounting}
For A10, let $C=2494$, $G_j=w_jg_j$, and
\[
Z(t)=\{j:r_j\le t<r_j+p_j\}.
\]
Let $q(t)$ be the minimum number of jobs removed from $Z(t)$ so that the
remaining GPU width is at most $C$. Then
\[
L_{\rm count}=\frac{1}{3600N}\int q(t)\,dt
\]
is a valid lower bound on ordinary mean waiting. The integral uses exact
half-open event intervals and integer GPU-width histograms. If $S(t)$ is the
actual service set, define $A(t)=S(t)\cap Z(t)$. A delayed job may still run
after $r_j+p_j$, so $S(t)\subseteq Z(t)$ is false in general. However,
$Z(t)\setminus A(t)$ is a deletion set leaving a feasible GPU total, so it has
at least $q(t)$ jobs. Therefore
\[
\sum_j\min(W_j,p_j)\ge\int q(t)\,dt,\qquad
\sum_jW_j\ge\int q(t)\,dt.
\]
The tuple $(r,p,s)=(0,2,4)$ illustrates the correction: at $t=5$ the job is
in $S(t)$ but not in $Z(t)$. Non-preemption gives the exact bookkeeping
identity $W_j=\min(W_j,p_j)+\max(W_j-p_j,0)$. The resulting three terms are
accounting around a relaxation, not causal shares or an exact distance to OPT.

\subsection{Complete A10 evidence}
The complete A10 input has 111,284 jobs and 2,494 nodes. The verified Q2 G2
SPT-HOL timing and the Q3 M0 SJF-HOL A10 timing have identical 111,284-row
ID, attribute, start, and finish records (canonical timing hash
\texttt{378acdbf...18baea8}). The full-cluster table below is read-only Q3
M0 FCFS/SJF-HOL evidence retrieved from fixed submission
\texttt{79e93c315f006d67f743b14a1d7faf16b00dd072}; its recorded run source is
commit \texttt{7075317883808b1cd7f80cadf3980be291c69cf0}. It is not a claim that the
Q2 engine ran a new 224,112-job trace.

\begin{table}[htbp]
\centering\small
\caption{Read-only Q3 M0 full-cluster summary, $H=90$ days; waiting uses complete drain.}
\begin{tabular}{llrrrrr}
\toprule
Policy & Pool & $n$ & Mean $W$ (h) & P95 (h) & Max (h) & $f\le H$\\
\midrule
FCFS & A10 &111284&900.129&1408.478&1423.254&54531\\
FCFS & A100-SXM4-80GB &80877&0.000&0.000&0.000&80101\\
FCFS & A800-SXM4-80GB &3959&0.146&0.000&12.250&3943\\
FCFS & GPU-series-1 &9643&0.029&0.000&9.971&8221\\
FCFS & GPU-series-2 &13633&0.045&0.000&6.992&13606\\
FCFS & H800 &4716&0.016&0.000&7.856&4705\\
FCFS & ALL &224112&446.971&1356.501&1423.254&165107\\
SJF-HOL & A10 &111284&32.906&11.087&3017.288&107002\\
SJF-HOL & A100-SXM4-80GB &80877&0.000&0.000&0.000&80101\\
SJF-HOL & A800-SXM4-80GB &3959&0.115&0.000&14.858&3943\\
SJF-HOL & GPU-series-1 &9643&0.434&0.000&76.818&8213\\
SJF-HOL & GPU-series-2 &13633&0.016&0.000&7.165&13606\\
SJF-HOL & H800 &4716&0.016&0.000&7.856&4705\\
SJF-HOL & ALL &224112&16.362&0.689&3017.288&217570\\
\bottomrule
\end{tabular}
\end{table}

For the A10 group, the job-count lower bound is
$8,941,062,908$ job-seconds, or $22.3179305700$ h. The Q1 fixed-pool
GPU-work-weighted lower bound is $560.5213678903$ h. With fixed
$H_j=w_jg_jp_j$ weights, the three persisted A10 timings give:

\begin{table}[htbp]
\centering
\caption{Ordinary and GPU-work-weighted A10 waiting; the two objectives need not rank policies alike.}
\begin{tabular}{lrrr}
\toprule
Policy & Mean job wait (h) & $\sum H_jW_j/\sum H_j$ (h) & Mean / job-count LB\\
\midrule
FCFS &900.129&619.352&40.33\\
SPT-HOL &32.906&677.690&1.47\\
SPT-FILL&34.967&677.496&1.57\\
\bottomrule
\end{tabular}
\end{table}

Thus SPT-HOL is the current ordinary-mean candidate, while the weighted
objective reverses the FCFS versus SPT ordering in this trace. Neither metric
alone establishes general optimality. The lower bound is conservative because
it relaxes CPU, node shape, and gang placement.

\subsection{Bounded protected-policy analysis and method choice}
On the verified 65,536-job A10 prefix, the ordinary means are 30.332 h
(SPT-HOL) and 31.042 h (SPT-FILL). URGENT-HOL/TIME/FIXED/FLEX are respectively
556.492/350.528/350.088/349.715 h; FIXED/FLEX accepted 62/133 crossings.
These are prefix protection-cost results only, not full-cluster or full-A10
urgent results. The FILL G3 cost record has 314,756 order builds. The
cross-review supplied 307,844 budget-exhausted events, giving 97.804% of
order-build records, but that ratio is a candidate-attempt diagnostic, not
97.804% of waiting time; the G3 archive does not independently persist that
event count.

For ordinary mean waiting, SPT-HOL is the current main candidate. SPT-FILL is
retained as a useful immediate-start baseline, and protected policies remain
valid safety or tail-preference experiments. SPT-HOL still has harmed jobs and
a slightly larger maximum wait than FILL; no dominance or general starvation
guarantee is claimed. All source hashes, row/ID checks, derived summaries,
and missing external full-cluster timing retention are recorded in
\texttt{results/20261016-b-q2-m2g-revision-01/}. The surrounding integrated
paper remains a draft; this file is the canonical Q2-only section.
